\section{执行摘要}

\begin{keyinsight}[结论先行]
\textbf{本报告是一份“如何研究”的证据备忘录，当前范围不含“现在买什么”的战略配置。} 全球机构的公开材料可以帮助个人识别经济引擎、制度前提与不可复制的能力；个人账户的资产配置比例、产品选择、预期收益、交易信号和执行权限均不由这些材料给出。当前结论是：所有市场敞口、策略袖套和候选对象均维持 \qzero{} 与 \noaction{}。
\end{keyinsight}

\section{本报告回答什么，不回答什么}

“战略配置”在机构语境中经常包含长期资本、跨资产风险预算、融资、流动性储备、衍生品对冲和治理机制。中国个人投资者即使观察到相同的宏观变量，也通常不拥有相同的账户权限、成本曲线、借券渠道、执行团队、数据权利或资产负债表。因此，本报告把问题收窄为：\textbf{哪些机构方法可以转化为公开信息下的研究纪律，哪些能力必须诚实地保留在机构边界之外。}

\begin{exhibitbox}[Exhibit 1：本备忘录的输出边界]
\small
\begin{tabularx}{\exhibitboxwidth}{L{3.2cm}X L{4.0cm}}
\toprule
\textbf{允许的读者用途} & \textbf{本报告提供的内容} & \textbf{明确不提供的内容} \\
\midrule
策略分类与研究学习 & 机构策略的经济引擎、数据/执行前提、可迁移与不可迁移的边界。 & “跟随某机构持仓”或把品牌当作买卖信号。 \\
纸面规则设计 & CTA、因子与公告事件的字段、时间戳、否决条件、纸面账本和复核流程。 & 实盘指令、外部订单、下单时点、账户操作或杠杆安排。 \\
制度与风险审查 & A 股、境内期货、产品和虚拟资产研究的约束清单。 & 法律意见、账户资格、适当性结论或跨境通道设计。 \\
反方检验 & 拥挤、模型失效、流动性、保证金、监管与行为风险的反例。 & 历史结果、模拟结果或机构业绩可以外推为个人未来回报。 \\
\bottomrule
\end{tabularx}
\sourcenote{S01--S15；本报告的研究边界设计。}
\end{exhibitbox}

\section{零数量作为证据状态}

\qzero{} 是一个硬门禁状态：在缺少完整而当前的市场、产品、合约、账户、风险预算和数据权利证据时，研究系统以零数量冻结，禁止生成任何非零数量、配置比例或外部行动。该状态只记录证据缺口与数量约束；市场方向和未来交易时点均不由此产生。零数量的作用是把下列常见跳跃拆开：

\begin{enumerate}
  \item 从“某机构公开谈论一种策略”跳到“个人可以复制该策略”；
  \item 从“一个历史回测或管理人业绩存在”跳到“个人未来能获得同类回报”；
  \item 从“某产品名称看起来相关”跳到“它具有同样经济暴露、流动性和结算条件”；
  \item 从“信号出现”跳到“可在真实市场、真实规则和真实成本下执行”。
\end{enumerate}

每一次跳跃都需要独立证据；更强烈的叙事不能替代证据。证据缺口存在时，最完整的研究结论可以是继续归零。

\section{读者应如何使用后文}

后文按“机构方法 $\rightarrow$ 可观察经济机制 $\rightarrow$ 中国市场边界 $\rightarrow$ 纸面研究规则 $\rightarrow$ 红队否决”展开。CTA、因子和事件章节列出具体参数，是为了使未来审查者能够复核规则是否被事后改写；这些参数属于\textbf{待验证的研究协议}，个人风险目标和仓位不在本报告输出范围内。所有数值都须与资料日期、数据版本、账户权限和成本假设同时审查，且在本报告中均不产生非零输出。

\begin{dashboardbox}[三条不可突破的底线]
\begin{tabularx}{\textwidth}{L{3.3cm}X}
\toprule
\textbf{底线} & \textbf{读者含义} \\
\midrule
不复制机构收益 & 管理人网页、AUM、媒体报道、历史 IRR、夏普比率或策略名称只能说明研究问题，不能构成个人预期回报。 \\
交易能力按证据逐项核验 & 稳定借券、主经纪商融资、场外衍生品、低延迟执行、做市库存、非公开信息和跨境账户均需单独证明，无法作为默认条件。 \\
不以规则文字代替事实核验 & 即使规则设计完整，任何具体产品、合约、日期和账户都须重新核验；若核验缺失，状态回到 \qzero{}。 \\
\bottomrule
\end{tabularx}
\end{dashboardbox}

\section{本报告与个人决策之间仍隔着什么}

即便一个研究协议本身逻辑清楚，个人层面的资产负债表、近期开支、现金流、税务居民身份、既有风险暴露、心理承受能力、账户权限和法律约束仍然不可由通用文本推断。此类信息及专业法律、税务、持牌机构和账户层面的意见均需独立提供与复核。本报告定位为：\textbf{可审计的研究框架；个人投资组合结论须另行形成。}

\begin{riskbox}[误用警报]
如果读者把“可迁移”理解为“应当交易”，把“纸面规则”理解为“已获批准”，或把“机构名称”理解为“收益背书”，则应停止使用本文并回到本章的 \qzero{} 状态。严谨的框架首先应允许自己不产生行动。
\end{riskbox}
